% questions.tex —— 2013 年·陕西卷·文科 · 题目内容片段（无答案）
% 由 tools/generate_projects.py 依来源总账自 2013/2013_数学真题_陕西文_试题.md 生成，勿手改（改源 MD 后重跑）


\begin{gknotice}
\textbf{来源说明}：本题面依据固定源 PDF 的两页 A3 题面逐题转录，原卷印刷页码为 839—840。第 839 面为选择题 1—10、填空题 11—17、解答题 18—19，第 840 面为解答题 20—23。第 15—17 题分别对应不等式选讲、几何证明选讲、坐标系与参数方程，按原卷分别编号并全部收录。

\textbf{解析归属}：固定源不附答案或解析；配套答案与解析均由本库按本题面独立复算。

\textbf{图形说明}：第 5 题频率分布直方图、第 12 题三视图、第 14 题内接矩形、第 16 题几何证明、第 20 题四棱柱，均由项目内 TikZ 矢量源重绘；第 17 题原页无配图。
\end{gknotice}

\gkpart{一、选择题}

\begin{question}
\gkti{1}设全集为 $\mathbf R$，函数 $f(x)=\sqrt{1-x}$ 的定义域为 $M$，则 $\complement_{\mathbf R}M$ 为\mbox{（　　）}

（A）$(-\infty,1)$；（B）$(1,+\infty)$；（C）$(-\infty,1]$；（D）$[1,+\infty)$。

\end{question}
\begin{question}
\gkti{2}已知向量 $\boldsymbol a=(1,m)$，$\boldsymbol b=(m,2)$，若 $\boldsymbol a\parallel\boldsymbol b$，则实数 $m$ 等于\mbox{（　　）}

（A）$-\sqrt2$；（B）$\sqrt2$；（C）$-\sqrt2$ 或 $\sqrt2$；（D）$0$。

\end{question}
\begin{question}
\gkti{3}设 $a$，$b$，$c$ 均为不等于 $1$ 的正实数，则下列等式中恒成立的是\mbox{（　　）}

（A）$\log_a b\cdot\log_c b=\log_c a$；（B）$\log_a b\cdot\log_c a=\log_c b$；（C）$\log_a(bc)=\log_a b\cdot\log_a c$；（D）$\log_a(b+c)=\log_a b+\log_a c$。

\end{question}
\begin{question}
\gkti{4}根据下列算法语句，当输入 $x$ 为 $60$ 时，输出 $y$ 的值为\mbox{（　　）}

\nopagebreak[4]
\[\begin{array}{|l|}
\hline
\text{\text{输入} }x;\\
\text{If }x\leqslant 50\text{ Then}\\
\qquad y=0.5*x\\
\text{Else}\\
\qquad y=25+0.6*(x-50)\\
\text{End If}\\
\text{\text{输出} }y.\\
\hline
\end{array}\]

（A）$25$；（B）$30$；（C）$31$；（D）$61$。

\Needspace{0.4\textheight}
\end{question}
\begin{question}
\gkti{5}对一批产品的长度（单位：毫米）进行抽样检测，下图为检测结果的频率分布直方图。根据标准，产品长度在区间 $[20,25)$ 上为一等品，在区间 $[15,20)$ 和 $[25,30)$ 上为二等品，在区间 $[10,15)$ 和 $[30,35]$ 上为三等品。用频率估计概率，现从该批产品中随机抽取 $1$ 件，则其为二等品的概率是\mbox{（　　）}
\nopagebreak[4]
\begin{center}% ============================================================
% f2013_shaanxi_arts_q05_histogram.tex
% 2013 陕西卷（文史类）第 5 题 频率分布直方图（TikZ 矢量重绘）
% 分组 [10,15)、[15,20)、[20,25)、[25,30)、[30,35]
% 纵轴为 频率/组距；柱高依次为 0.02、0.04、0.06、0.05、0.03
% 横轴 0 与 10 之间有断裂记号
% ============================================================
\begin{tikzpicture}[
  x=1cm, y=1cm, line width=0.5pt, font=\footnotesize, >=Stealth
]

  % ---------- 坐标轴 ----------
  \draw[->] (0,0) -- (7.30,0);
  \draw[->] (0,0) -- (0,4.95);

  % 横轴原点附近的断裂记号
  \draw (0.16,0.09) -- (0.28,-0.09);
  \draw (0.34,0.09) -- (0.46,-0.09);

  % ---------- 频率/组距 纵轴标题 ----------
  \node at (-0.78,4.50) {$\dfrac{\text{频率}}{\text{组距}}$};

  % ---------- 五根柱 ----------
  \draw (0.80,0) rectangle (2.00,1.20);   % [10,15)
  \draw (2.00,0) rectangle (3.20,2.40);   % [15,20)
  \draw (3.20,0) rectangle (4.40,3.60);   % [20,25)
  \draw (4.40,0) rectangle (5.60,3.00);   % [25,30)
  \draw (5.60,0) rectangle (6.80,1.80);   % [30,35]

  % ---------- 纵轴刻度与虚线 ----------
  \foreach \y/\lab in {1.20/0.02, 1.80/0.03, 2.40/0.04, 3.60/0.06}{
    \draw (-0.07,\y) -- (0,\y);
    \node[left] at (-0.10,\y) {\lab};
    \draw[dashed, gray!55] (0,\y) -- (6.80,\y);
  }

  % ---------- 横轴刻度 ----------
  \foreach \x/\lab in {0/0, 0.80/10, 2.00/15, 3.20/20, 4.40/25, 5.60/30, 6.80/35}{
    \draw (\x,0) -- (\x,-0.08);
    \node[below] at (\x,-0.10) {\lab};
  }

  % ---------- 横轴标题 ----------
  \node[anchor=west] at (6.95,0.26) {长度/毫米};

\end{tikzpicture}
\end{center}
\nopagebreak[4]
（A）$0.09$；（B）$0.20$；（C）$0.25$；（D）$0.45$。

\end{question}
\begin{question}
\gkti{6}设 $z$ 是复数，则下列命题中的假命题是\mbox{（　　）}

（A）若 $z^2\geqslant0$，则 $z$ 是实数；（B）若 $z^2<0$，则 $z$ 是虚数；（C）若 $z$ 是虚数，则 $z^2\geqslant0$；（D）若 $z$ 是纯虚数，则 $z^2<0$。

\end{question}
\begin{question}
\gkti{7}若点 $(x,y)$ 位于曲线 $y=|x|$ 与 $y=2$ 所围成的封闭区域，则 $2x-y$ 的最小值是\mbox{（　　）}

（A）$-6$；（B）$-2$；（C）$0$；（D）$2$。

\end{question}
\begin{question}
\gkti{8}已知点 $M(a,b)$ 在圆 $O:x^2+y^2=1$ 外，则直线 $ax+by=1$ 与圆 $O$ 的位置关系是\mbox{（　　）}

（A）相切；（B）相交；（C）相离；（D）不确定。

\end{question}
\begin{question}
\gkti{9}设 $\triangle ABC$ 的内角 $A$，$B$，$C$ 所对的边分别为 $a$，$b$，$c$，若 $b\cos C+c\cos B=a\sin A$，则 $\triangle ABC$ 的形状为\mbox{（　　）}

（A）直角三角形；（B）锐角三角形；（C）钝角三角形；（D）不确定。

\end{question}
\begin{question}
\gkti{10}设 $[x]$ 表示不大于 $x$ 的最大整数，则对任意实数 $x$，有\mbox{（　　）}

（A）$[-x]=-[x]$；（B）$\left[x+\dfrac12\right]=[x]$；（C）$[2x]=2[x]$；（D）$[x]+\left[x+\dfrac12\right]=[2x]$。

\end{question}
\begingroup
\renewcommand{\gkpart}[1]{\Needspace{0.2\textheight}\subsection*{#1}\nopagebreak}
\gkpart{二、填空题}
\endgroup

\begin{question}
\gkti{11}双曲线 $\dfrac{x^2}{16}-\dfrac{y^2}{9}=1$ 的离心率为\gkblank{}。

\Needspace{0.4\textheight}
\end{question}
\begin{question}
\gkti{12}某几何体的三视图如图所示，则其表面积为\gkblank{}。
\nopagebreak[4]
\begin{center}% ============================================================
% f2013_shaanxi_arts_q12_views.tex
% 2013 陕西卷（文史类）第 12 题 三视图（TikZ 矢量重绘）
% 主视图、左视图：直径 2、高 1 的半圆；俯视图：直径 2 的圆
% 由此还原为半径 1 的半球（表面积 3π）
% ============================================================
\begin{tikzpicture}[
  x=1cm, y=1cm, line width=0.5pt, font=\footnotesize, >=Stealth
]

  % ---------- 主视图 ----------
  \draw (0,0) -- (2,0);
  \draw (0,0) arc (180:0:1);
  \draw[|<->|] (-0.48,0) -- (-0.48,1) node[midway, left] {$1$};
  \foreach \y in {0,1} \draw (-0.33,\y) -- (-0.63,\y);
  \draw[|<->|] (0,-0.55) -- (2,-0.55) node[midway, below] {$2$};
  \foreach \x in {0,2} \draw (\x,-0.70) -- (\x,-0.40);
  \node at (1,-1.18) {主视图};

  % ---------- 左视图 ----------
  \begin{scope}[xshift=3.60cm]
    \draw (0,0) -- (2,0);
    \draw (0,0) arc (180:0:1);
    \draw[|<->|] (0,-0.55) -- (2,-0.55) node[midway, below] {$2$};
    \foreach \x in {0,2} \draw (\x,-0.70) -- (\x,-0.40);
    \node at (1,-1.18) {左视图};
  \end{scope}

  % ---------- 俯视图 ----------
  \draw (1,-3.55) circle (1);
  \node at (1,-4.98) {俯视图};

\end{tikzpicture}
\end{center}
\nopagebreak[4]
\end{question}
\begin{question}
\gkti{13}观察下列等式：
\nopagebreak[4]
\[(1+1)=2\times1,\]
\nopagebreak[4]
\[(2+1)(2+2)=2^2\times1\times3,\]
\nopagebreak[4]
\[(3+1)(3+2)(3+3)=2^3\times1\times3\times5,\]
\nopagebreak[4]
\[\cdots\cdots\]
照此规律，第 $n$ 个等式可为\gkblank{}。

\Needspace{0.4\textheight}
\end{question}
\begin{question}
\gkti{14}在如图所示的锐角三角形空地中，欲建一个面积最大的内接矩形花园（阴影部分），则其边长 $x$ 为\gkblank{}（m）。
\nopagebreak[4]
\begin{center}% ============================================================
% f2013_shaanxi_arts_q14_triangle.tex
% 2013 陕西卷（文史类）第 14 题 锐角三角形空地内接矩形（TikZ 矢量重绘）
% 底 40 m、高 40 m，内接矩形一边落在底边上，宽为 x
% 示意图不按原比例；矩形宽高关系满足 x/40 + y/40 = 1
% ============================================================
\begin{tikzpicture}[
  x=1cm, y=1cm, line width=0.5pt, font=\footnotesize, >=Stealth
]

  % ---------- 三角形空地 ----------
  \draw (0,0) -- (4,0) -- (2,4) -- cycle;

  % ---------- 内接矩形（阴影） ----------
  \fill[gray!25] (1.20,0) rectangle (2.80,2.40);
  \draw (1.20,0) rectangle (2.80,2.40);
  \node at (2,2.62) {$x$};

  % ---------- 高的标注（40 m） ----------
  \draw[dashed, gray!70] (2,4) -- (4.75,4);
  \draw[dashed, gray!70] (4,0) -- (4.75,0);
  \draw[|<->|] (4.55,0) -- (4.55,4) node[midway, right] {$40$ m};
  \foreach \y in {0,4} \draw (4.38,\y) -- (4.72,\y);

  % ---------- 底的标注（40 m） ----------
  \draw[|<->|] (0,-0.58) -- (4,-0.58) node[midway, below] {$40$ m};
  \foreach \x in {0,4} \draw (\x,-0.73) -- (\x,-0.43);

\end{tikzpicture}
\end{center}
\nopagebreak[4]
\end{question}
\begin{question}
\gkti{15}设 $a$，$b\in\mathbf R$，$|a-b|>2$，则关于实数 $x$ 的不等式 $|x-a|+|x-b|>2$ 的解集是\gkblank{}。

\Needspace{0.4\textheight}
\end{question}
\begin{question}
\gkti{16}如图，$AB$ 与 $CD$ 相交于点 $E$，过 $E$ 作 $BC$ 的平行线与 $AD$ 的延长线交于点 $P$，已知 $\angle A=\angle C$，$PD=2DA=2$，则 $PE=$\gkblank{}。
\nopagebreak[4]
\begin{center}% ============================================================
% f2013_shaanxi_arts_q16_geometry.tex
% 2013 陕西卷（文史类）第 16 题 圆与平行线几何证明（TikZ 矢量重绘）
% P、D、A 共线，PD=2、DA=1；AB 与 CD 交于 E；PE ∥ BC
% 图中 PE 与 BC 严格平行；∠A=∠C 由作图满足（示意图不按原比例）
% ============================================================
\begin{tikzpicture}[
  x=1cm, y=1cm, line width=0.5pt, font=\footnotesize, >=Stealth
]

  % ---------- 主要线段 ----------
  \draw (-3.400,0) -- (1.700,0);              % P--D--A
  \draw (1.700,0) -- (-1.745,1.730);           % AB
  \draw (3.576,2.353) -- (0,0);                % CD
  \draw (-3.400,0) -- (0.736,0.484);           % PE
  \draw (-1.745,1.730) -- (3.576,2.353);       % BC

  % ---------- 顶点标注 ----------
  \node[below left] at (-3.400,0) {$P$};
  \node[below] at (0,0) {$D$};
  \node[below right] at (1.700,0) {$A$};
  \node[below right] at (0.736,0.484) {$E$};
  \node[above left] at (-1.745,1.730) {$B$};
  \node[above right] at (3.576,2.353) {$C$};

\end{tikzpicture}
\end{center}
\nopagebreak[4]
\end{question}
\begin{question}
\gkti{17}圆锥曲线 $\begin{cases}x=t^2,\\y=2t,\end{cases}$（$t$ 为参数）的焦点坐标是\gkblank{}。

\end{question}
\begingroup
\renewcommand{\gkpart}[1]{\Needspace{0.2\textheight}\subsection*{#1}\nopagebreak}
\gkpart{三、解答题}
\endgroup

\begin{question}
\gkti{18}已知向量 $\boldsymbol a=\left(\cos x,-\dfrac12\right)$，$\boldsymbol b=\left(\sqrt3\sin x,\cos2x\right)$，$x\in\mathbf R$，设函数 $f(x)=\boldsymbol a\cdot\boldsymbol b$。

（1） 求 $f(x)$ 的最小正周期。

（2） 求 $f(x)$ 在 $\left[0,\dfrac\pi2\right]$ 上的最大值和最小值。

\end{question}
\begin{question}
\gkti{19}设 $S_n$ 表示数列 $\{a_n\}$ 的前 $n$ 项和。

（1） 若 $\{a_n\}$ 是等差数列，推导 $S_n$ 的计算公式；

（2） 若 $a_1=1$，$q\ne0$，且对所有正整数 $n$，有 $S_n=\dfrac{1-q^n}{1-q}$，判断 $\{a_n\}$ 是否为等比数列，并证明你的结论。

\Needspace{0.4\textheight}
\end{question}
\begin{question}
\gkti{20}如图，四棱柱 $ABCD-A_1B_1C_1D_1$ 的底面 $ABCD$ 是正方形，$O$ 为底面中心，$A_1O\perp$ 平面 $ABCD$，$AB=AA_1=\sqrt2$。

（1） 证明：平面 $A_1BD\parallel$ 平面 $CD_1B_1$；

（2） 求三棱柱 $ABD-A_1B_1D_1$ 的体积。
\nopagebreak[4]
\begin{center}% ============================================================
% f2013_shaanxi_arts_q20_prism.tex
% 2013 陕西卷（文史类）第 20 题 四棱柱（TikZ 矢量重绘）
% 底面 ABCD 为正方形（斜二测示意），O 为底面中心，A1O ⊥ 底面
% 虚线：被遮挡的底边、侧棱、两平面的对角线以及 A1O
% ============================================================
\begin{tikzpicture}[
  x=1cm, y=1cm, line width=0.5pt, font=\footnotesize, >=Stealth
]

  % ---------- 底面（正方形示意） ----------
  \draw (0,0) -- (3,0) -- (4.6,1.2);
  \draw[dashed] (4.6,1.2) -- (1.6,1.2) -- (0,0);

  % ---------- 侧棱 ----------
  \draw (0,0) -- (2.3,3.0);         % AA1
  \draw (3,0) -- (5.3,3.0);         % BB1
  \draw (4.6,1.2) -- (6.9,4.2);     % CC1
  \draw[dashed] (1.6,1.2) -- (3.9,4.2);   % DD1

  % ---------- 上底面 ----------
  \draw (2.3,3.0) -- (5.3,3.0) -- (6.9,4.2) -- (3.9,4.2) -- cycle;

  % ---------- 平面 A1BD 与 平面 CD1B1 的对角线 ----------
  \draw[dashed] (2.3,3.0) -- (3,0);
  \draw[dashed] (2.3,3.0) -- (1.6,1.2);
  \draw[dashed] (3,0) -- (1.6,1.2);
  \draw[dashed] (4.6,1.2) -- (3.9,4.2);
  \draw[dashed] (4.6,1.2) -- (5.3,3.0);
  \draw[dashed] (3.9,4.2) -- (5.3,3.0);

  % ---------- 高 A1O ----------
  \draw[dashed] (2.3,3.0) -- (2.3,0.6);

  % ---------- 顶点标注 ----------
  \node[below left]  at (0,0) {$A$};
  \node[below right] at (3,0) {$B$};
  \node[below right] at (4.62,1.12) {$C$};
  \node[left]        at (1.30,1.45) {$D$};
  \node[left]        at (2.16,3.04) {$A_1$};
  \node[above]       at (5.3,3.08) {$B_1$};
  \node[above right] at (6.92,4.22) {$C_1$};
  \node[above]       at (3.9,4.24) {$D_1$};
  \node[below]       at (2.12,0.42) {$O$};

\end{tikzpicture}
\end{center}
\nopagebreak[4]
\end{question}
\begin{question}
\gkti{21}有 $7$ 位歌手（$1$ 至 $7$ 号）参加一场歌唱比赛，由 $500$ 名大众评委现场投票决定歌手名次。根据年龄将大众评委分为五组，各组的人数如下：

\begin{center}
\begin{tabular}{llllll}
\toprule
组别 & $A$ & $B$ & $C$ & $D$ & $E$ \\
\midrule
人数 & $50$ & $100$ & $150$ & $150$ & $50$ \\
\bottomrule
\end{tabular}
\end{center}

（1） 为了调查评委对 $7$ 位歌手的支持情况，现用分层抽样方法从各组中抽取若干评委，其中从 $B$ 组抽取了 $6$ 人，请将其余各组抽取的人数填入下表。

\begin{center}
\begin{tabular}{llllll}
\toprule
组别 & $A$ & $B$ & $C$ & $D$ & $E$ \\
\midrule
人数 & $50$ & $100$ & $150$ & $150$ & $50$ \\
抽取人数 &  & $6$ &  &  &  \\
\bottomrule
\end{tabular}
\end{center}

（2） 在（1）中，若 $A$，$B$ 两组被抽到的评委中各有 $2$ 人支持 $1$ 号歌手，现从这两组被抽到的评委中分别任选 $1$ 人，求这 $2$ 人都支持 $1$ 号歌手的概率。

\end{question}
\begin{question}
\gkti{22}已知动点 $M(x,y)$ 到直线 $l:x=4$ 的距离是它到点 $N(1,0)$ 的距离的 $2$ 倍。

（1） 求动点 $M$ 的轨迹 $C$ 的方程；

（2） 过点 $P(0,3)$ 的直线 $m$ 与轨迹 $C$ 交于 $A$，$B$ 两点，若 $A$ 是 $PB$ 的中点，求直线 $m$ 的斜率。

\end{question}
\begin{question}
\gkti{23}已知函数 $f(x)=\mathrm e^x$，$x\in\mathbf R$。

（1） 求 $f(x)$ 的反函数的图象上点 $(1,0)$ 处的切线方程；

（2） 证明：曲线 $y=f(x)$ 与曲线 $y=\dfrac12x^2+x+1$ 有唯一公共点；

（3） 设 $a<b$，比较 $f\left(\dfrac{a+b}{2}\right)$ 与 $\dfrac{f(b)-f(a)}{b-a}$ 的大小，并说明理由。

\end{question}